% Generated draft: requires mathematical and source review.
\documentclass[11pt]{article}
\usepackage[margin=25mm]{geometry}
\usepackage{fontspec}
\setmainfont{Latin Modern Roman}
\usepackage{amsmath,amssymb,mathtools,unicode-math}
\setmathfont{Latin Modern Math}
\usepackage{enumitem,needspace,xurl,hyperref}
\setlength{\parindent}{0pt}
\setlength{\parskip}{5pt}
\setlength{\emergencystretch}{3em}
\newcommand{\N}{\mathbb{N}}
\newcommand{\Z}{\mathbb{Z}}
\newcommand{\Q}{\mathbb{Q}}
\newcommand{\R}{\mathbb{R}}
\newcommand{\C}{\mathbb{C}}
\newcommand{\F}{\mathbb{F}}
\newcommand{\A}{\mathbb{A}}
\newcommand{\PP}{\mathbb{P}}
\newcommand{\E}{\mathbb{E}}
\newcommand{\eps}{\varepsilon}
\newcommand{\dd}{\,\mathrm{d}}
\begin{document}
\section*{UnsolvedMath ID: 10600048}
\textbf{Source problem code:} catalogue entry
\subsection{Definitions and mathematical statement}
Let $\mathcal V$ and $\mathcal C$ denote virtual and classical knot types, and let $\iota:\mathcal C\to\mathcal V$ be the usual inclusion. Problem-list item 48 asks for a functorial construction projecting virtual knots to classical knots, analogous to the source's previously constructed map from virtual knots to virtual knots with orientable atoms.

If ``projection'' means a retraction, the requested map $P:\mathcal V\to\mathcal C$ should satisfy
\[
 P\circ\iota=\operatorname{id}_{\mathcal C},
\]
and be induced by a diagrammatic construction compatible with the relevant moves. The excerpt does not explicitly give this retraction identity or specify the morphisms for functoriality, so this is the natural conditional reading rather than a completed definition of the source's functor. A mere constant map to a classical knot would not express the intended projection.

The source proposes using an index defined through characteristic classes and says that a condition arising under the third Reidemeister move should hold, but supplies neither the index nor that condition. It also observes that for homologically trivial knots a two-branched covering can be considered as a projection; the ambient surface, homology coefficients, branch data, and resulting map are omitted. These suggestions remain part of the question, not established constructions in this draft.

The registry reports that a functorial map to classical knots was constructed in 2010, with further work in 2023. Consequently the original existence request and its later reported progress must be distinguished; no precise residual structural problem is stated in the supplied summary.
\subsection{Short English statement}
Construct a move-compatible projection from virtual to classical knots, analogous to the orientable-atom construction. The registry reports a later construction, but the exact functorial conditions and remaining question need source verification.
\subsection{Sources}
\begin{itemize}
\item Local registry: \url{data/unsolvedmath/problems.json}, permanent UnsolvedMath ID $10600048$, exact source code \texttt{catalogue entry}.
\item \emph{Unsolved Problems in Virtual Knot Theory and Combinatorial Knot Theory}, revised survey, problem-list item 48: \url{https://arxiv.org/abs/1409.2823}. Supplied access date: 2026-07-29; extraction: source-tex; compound item retained intact.
\item Supplied progress references: ``A Functorial Map from Virtual Knots to Classical Knots and Generalisations of Parity,'' \texttt{arXiv:1011.xxxx} (2010), and ``Local transformations and functorial maps,'' \texttt{arXiv:2301.xxxx} (2023). Both identifiers are incomplete placeholders; no authors or exact references are supplied here.
\end{itemize}
\end{document}
